% META: {"id":"TNSI-LANG-EX-004","chapitre":"TNSI-LANGAGES-ET-PROGRAMMATION","type_objet":"exercice","capacites":["T-LANG-05"],"capacites_codes":["C11"],"methodes":[],"parcours":1,"competences":["analyser"],"duree_min":15,"mode_creation":"ex_nihilo","sources_inspiration":[],"fichier_tex":"chapitres/TNSI-LANGAGES-ET-PROGRAMMATION/exercices/TNSI-LANG-EX-004.tex","corrige_tex":"chapitres/TNSI-LANGAGES-ET-PROGRAMMATION/corriges/TNSI-LANG-CO-004.tex","status":"approved","origin":"generated","status_history":["generated","audited","accepted","approved"],"release_acceptance":"RELEASE_OWNER_BATCH_ACCEPTANCE","acceptance_closure_digest":"sha256:9b3ccf9a81c5520fbb7e03b7b2d3e7bf2057a02834b6d908bf3be5f49f3b553f","acceptance_receipt_digest":"sha256:4fad86f0282e0fa4e0419cca1ad6379ae7af5a280c04a4a90880f90a1c044242"}
% BEGIN-VERIFY
% def ajouter_note(bulletin, note):
%     bulletin.append(note)
%     return bulletin
% bulletin_zero = []
% b1 = ajouter_note(bulletin_zero, 15)
% b2 = ajouter_note(bulletin_zero, 12)
% assert b1 is b2
% assert b1 == [15, 12]
% def ajouter_note_corrige(bulletin, note):
%     return bulletin + [note]
% bulletin_zero2 = []
% b1c = ajouter_note_corrige(bulletin_zero2, 15)
% b2c = ajouter_note_corrige(bulletin_zero2, 12)
% assert b1c == [15]
% assert b2c == [12]
% assert bulletin_zero2 == []
% END-VERIFY
\begin{exercice}{TNSI-LANG-EX-004}{1}{15}
Le code suivant contient un bug lié à un \textbf{effet de bord} inattendu :
\begin{python}
def ajouter_note(bulletin, note):
    bulletin.append(note)
    return bulletin

bulletin_A = []
bulletin_B = ajouter_note(bulletin_A, 15)
bulletin_C = ajouter_note(bulletin_A, 12)
\end{python}
\begin{enumerate}
  \item Que contiennent \lstinline{bulletin_A}, \lstinline{bulletin_B} et \lstinline{bulletin_C} après ces trois lignes ? Le résultat est-il celui attendu si l'on voulait deux bulletins \textbf{indépendants} (un avec la note $15$, un autre avec la note $12$) ?
  \item Corriger \lstinline{ajouter_note} pour qu'elle renvoie un \textbf{nouveau} bulletin sans modifier celui passé en argument.
\end{enumerate}
\end{exercice}
